\scanpage{131}
$V$ anneau de valuation discrète, corps résiduel $k$ parfait, corps des fractions $K$.

Pour tout préschéma $X$ sur $V$, on définit $\bar X$ par Greenberg.
\[
\bar X(k')\simeq X(V_{k'})\deleted{\unsure{\longrightarrow}X}
\]
\margin{($k'$ \unsure{alg.} sur $k$, $V_{k'}$ \unsure{ext. correspondante de} $V$)}

Si $k'$ est parfait, $V_{k'}$ est un anneau de valuation discrète, soit $K_{k'}=K'$ son corps des fractions, on a
\[
\bar X(\unsure{k|k'})\simeq X(V_{k'})\longrightarrow X(K')
\]

Soit $\deleted{\unsure{\varphi:X\longrightarrow}}Y$ un autre $V$ préschéma, \inserted{\uline{propre sur} $V$,} et soit
\[
\varphi:X_K\longrightarrow Y_K
\]
un morphisme de $K$-préschémas, nous allons en déduire une \uline{application ensembliste}
\[
\bar\varphi:\bar X\longrightarrow\bar Y.
\]

Soit en effet $\bar x\in\bar X$, soit $k'=\widetilde{k}(\bar x)$ \deleted{une} \inserted{``la''} clôture parfaite de $k(\bar x)$, alors $\bar x$ définit un p\textsuperscript{t} $\bar x'\in\bar X(k')$, d'où un p\textsuperscript{t} \deleted{\unsure{générique}} $\alpha_X(\bar x')\in\bar X(V_{k'})\subset\bar X(K_{k'})$ et appliquant un point
\[
\varphi(\alpha_X(\bar x'))\in\overset{\deleted{\rule{.6em}{.3pt}}}{Y}(K_{k'})
\xleftarrow{\sim}\overset{\deleted{\rule{.6em}{.3pt}}}{Y}(V_{k'})
\xleftarrow[\sim]{\alpha_Y}\bar Y(k'),
\]
dont la \unsure{localité} est notée $\bar\varphi(x)$. On a donc par construction un homo
\[
k(\varphi(\bar x))\longrightarrow\widetilde{k(\bar x)}
\]
(ici le $\sim$ dénote la clôture parfaite).

\margin{Si $\overset{\deleted{\rule{.6em}{.3pt}}}{Y}$ n'est pas propre sur $V$ et $\varphi$ seulement \unsure{application rationnelle}, on aura une ``application \unsure{non} \unsure{partout définie}'' $\bar X\longrightarrow\bar Y$, ayant les propriétés}

Par la suite, $X,Y$ de type fini sur $K$, et \inserted{\unsure{géom.}} intègres

\uline{Proposition 1} Supposons $\bar x\in\bar X$ \uline{épais} et $\varphi$ \uline{dominant gén.\textsuperscript{t} séparable} [i.e. fibre gén. séparable]. Alors \inserted{\unsure{si|1\textsuperscript{o}}} $\bar y=\varphi(\bar x)$ \inserted{\unsure{existe, et|étale, et}} est \uline{épais}, $k(\bar y)\subset k(\bar x)$, et l'application ponctuelle $\bar x_{\mathrm{red}}\longrightarrow\bar y_{\mathrm{red}}$ définie par cette inclusion provient d'un \unsure{promorphisme} d'indice fini sur un ouvert constructible dense $U\in\bar x$, [d'image un ouvert constructible de $\bar y_{\mathrm{red}}$].
\figure{\begin{tikzpicture}\draw (0,0)--(0,-2);\end{tikzpicture}}
\margin{130}
